% 原文 pp.218--222
\subsubsectionlink{スピン・ボソン模型のBorn--Markovマスター方程式}
前の5.3.1節で論じた、簡単化したスピン・ボソン模型では、全Hamiltonian(5.90)は中心スピンの$\hat\sigma_z$固有基底で完全に対角的だった。これは、スピン系が$\{|0\rangle,|1\rangle\}$基底で固有の運動をもたないことを意味する。系が環境に結合していなければ、測定で系を$|0\rangle$または$|1\rangle$に見いだす確率は、時間とともに変化しない。

次に、Hamiltonian(5.90)へトンネル項$-\Delta_0\hat\sigma_x/2$を加え、固有の運動を含むように、この模型を一般化しよう。$\Delta_0$はトンネル行列要素である［(5.1)参照］。環境座標を$(\hat q_i,\hat p_i)$表示に戻すと、一般的なスピン・ボソン模型はHamiltonian
\chfiveeq{136}
で定義される。Leggettら\cite{ref209}は、このHamiltonianの形を物理的考察から厳密に動機づける方法を、詳しく示している。ここでの目的には、スピン系と振動子環境の模型の役割についての5.1節の導入的な説明で、Hamiltonian(5.136)への関心を動機づけるには十分であるとする。

このHamiltonianは、5.3.1節で考えたトンネル効果のないHamiltonian(5.90)より、わずかに複雑になっただけのように見える。それにもかかわらず、トンネル項の存在によって、(5.136)で定義する一般的スピン・ボソン模型は、数学的解法の点でも、得られるデコヒーレンスの動力学の、パラメータの相対的大きさへの依存性の点でも、はるかに複雑になる。前述のLeggettらの総説\cite{ref209}は、この模型の厳密解と近似解を、非常に詳しく論じている。

ここでは、スピン・ボソン模型のBorn--Markovマスター方程式の導出に注目する。系が大きな環境に弱く結合する場合、このマスター方程式が適切であることを思い出そう。以下の計算を容易にするため、Hamiltonian(5.136)に、小さな簡単化を一つ導入する。

% 原文 p.219
\hypertarget{srcpage:219}{}\label{srcpage:219}
すなわち、中心系が対称な二重井戸ポテンシャルを運動する場合に対応して、非対称エネルギー$\omega_0$がゼロであると仮定する。この仮定により、(5.136)の最初の$\hat\sigma_z$の項が消え、模型はHamiltonian
\chfiveeq{137}
で定義される。この、(5.136)をわずかに簡単化したHamiltonianの形によって、Born--Markovマスター方程式の導出が見通しよくなり、量子Brown運動との興味深い類似性も、より明確に確立できる。

幸い、導出は量子Brown運動の場合（5.2.1節参照）と非常に似た方法で進められる。実際、同じ種類の環境と、双線形の系・環境結合を扱うので、多くの結果をそのまま引き継げる。例えば、量子Brown運動のマスター方程式の導出の大部分は、環境自己相関関数(5.15)の計算に費やされた。これらの関数は中心系の表現に依存しないので、以下では、5.2.1節で得た結果を直接用いられる。

特に、量子Brown運動のマスター方程式の形(5.22)を引き継ぎ、位置演算子$\hat X$をPauliの$z$スピン演算子$\hat\sigma_z$に置き換えるだけでよい。
\chfiveeq{138}
ここで$\hat H_{\mathrm S}=-\Delta_0\hat\sigma_x/2$は系のHamiltonianである。(5.138)の雑音核$\nu(\tau)$と散逸核$\eta(\tau)$は、量子Brown運動の場合とまったく同じ形(5.16)、(5.17)、すなわち
\chfiveeq{139_140}
を取る。相互作用描像での演算子$\hat\sigma_z$の時間依存性は、系のHamiltonian $\hat H_{\mathrm S}$についてのHeisenberg運動方程式を解けば計算できる。結果は
\chfiveeq{141}
である。

% 原文 p.220
\hypertarget{srcpage:220}{}\label{srcpage:220}
この表式が、量子Brown運動模型で用いた相互作用描像の位置演算子$\hat X(\tau)$の結果(5.24)と、形式的に似ていることに注意しよう。

(5.141)をマスター方程式(5.138)へ代入する。項を整理し、Pauliスピン演算子の標準的恒等式$\hat\sigma_y\hat\sigma_z=\mathrm i\hat\sigma_x$、$\hat\sigma_z\hat\sigma_y=-\mathrm i\hat\sigma_x$を用いる、直接的な代数計算の後、Hamiltonian(5.137)で定義するスピン・ボソン模型のBorn--Markovマスター方程式の最終形
\chfiveeq{142}
を得る。

この方程式には、いくつかの新しい量が現れる。まず
\chfiveeq{143}
は、系の繰り込まれた、一般には非HermiteのHamiltonianである。すなわち、トンネル項$-\Delta_0/2$は、環境の影響によって$\zeta^*$だけ変わり、
\chfiveeq{144}
である。(5.142)、(5.144)に現れる係数$\widetilde D$、$\widetilde f$、$\widetilde\gamma$は、量子Brown運動模型の通常拡散係数$D$、異常拡散係数$f$［(5.26c)、(5.26d)参照］、減衰係数$\gamma$［(5.26b)参照］に、それぞれ類似した表式
\chfiveeq{145a_145b_145c}
で決まる。したがって、環境のスペクトル密度の具体形を選べば、係数(5.145)は、量子Brown運動の場合（5.2.3節参照）とほぼ同じ方法で評価できる。

マスター方程式(5.142)の右辺の各項の解釈も、量子Brown運動の場合と似ている。項
\chfiveeq{146}
は、標準のLindblad二重交換子の形(4.45)である。

% 原文 p.221
\hypertarget{srcpage:221}{}\label{srcpage:221}
相互作用Hamiltonianの構造［(5.136)参照］からすでに明らかなように、これは、環境による系の観測量$\hat\sigma_z$の直接監視を表す。したがって、係数$\widetilde D$で与えられる率での、基底$\{|0\rangle,|1\rangle\}$のデコヒーレンスを記述する。

このことは、(5.146)を$\{|0\rangle,|1\rangle\}$基底の$2\times2$行列として表しても、具体的に見られる。縮約密度行列の要素$\langle i|\hat\rho_{\mathrm S}(t)|j\rangle$、$i\in\{0,1\}$を$\rho_{\mathrm S}^{(ij)}(t)$と書けば、
\chfiveeq{147}
を得る。したがって、マスター方程式の他のすべての項を無視し、$\widetilde D$の項(5.146)だけによる、$\hat\sigma_z$固有基底での縮約密度行列の非対角要素の時間発展は
\chfiveeq{148}
である。これは、率$\widetilde D$での非対角要素の指数関数的減衰を意味する。対角要素$\rho_{\mathrm S}^{(00)}(t)$、$\rho_{\mathrm S}^{(11)}(t)$、すなわち占有確率は、項(5.146)が生成する運動の影響を受けず、減衰を伴わない純粋なデコヒーレンス過程であることを示す。このことも、通常拡散項が空間的コヒーレンスを指数関数的に減衰させた、量子Brown運動の場合［(5.35)参照］に類似している。

トンネル行列要素が消える$\Delta_0=0$の場合、$\zeta=0$でもあることに注意しよう。このときマスター方程式(5.142)は、減衰や振動数シフトに対応する項を含まない「純粋デコヒーレンス」の形
\chfiveeq{149}
に帰着する。これは、Lindblad演算子$\hat L=\hat\sigma_z$をもつLindblad方程式(4.45)にすぎず、環境による系の$z$スピン値の監視を記述する。

最後に、マスター方程式(5.142)の右辺の最後の二項は、係数$\widetilde f$、$\widetilde\gamma$［(5.145b)、(5.145c)参照］の兼ね合いで定量化される、二準位系の崩壊を記述する。これは、例えば各項を行列形で書き下せば、かなり容易にわかる（ここではデコヒーレンスに注目するため、証明は示さない）。

% 原文 p.222
\hypertarget{srcpage:222}{}\label{srcpage:222}
まとめると、スピン・ボソン模型のBorn--Markovマスター方程式は、量子Brown運動のマスター方程式の導出と、非常に類似した方法で得られることがわかった。次の段階は、Born--Markov近似を超えて、スピン・ボソン模型を調べることである。この近似は、中心二準位系と環境が非常に低温にある物理的状況では、しばしば破れる（4.4節の議論参照）。そのような低温系の顕著な例は、絶対零度に近い温度で動作する超伝導量子ビット系である（6.3節参照）。

前述のように、スピン・ボソン模型の豊かな非Markov動力学は、Leggettら\cite{ref209}が非常に詳しく解析している。予想どおり、その運動は、環境温度、スペクトル密度の形（subohmic、Ohmic、supraohmic）、系と環境の全体的結合強度など、模型のさまざまなパラメータに強く依存する。それぞれのパラメータ領域には、局在化、指数関数的または非コヒーレントな緩和、指数関数的減衰、強くまたは弱く減衰したコヒーレント振動という、特徴的な運動が現れる（パラメータ領域と対応する動力学の概観は、文献\cite{ref209}の表I参照）。

ここでは、このかなり複雑な動力学は論じず、上で調べた二つの模型、すなわち簡単化したスピン・ボソン模型と、Born--Markov近似の枠組みでのスピン・ボソン模型で満足することにしよう。これらの模型は、スピン・ボソン型模型のデコヒーレンス研究に関係する主要な特徴を、すでに示している。スピン・ボソン模型の散逸とデコヒーレンスの非Markov動力学を深く解析したい読者は、Leggettらの論文\cite{ref209}とWeissの著書\cite{ref218}に、膨大な情報を見いだすだろう。
